\documentclass[12pt]{article}
\usepackage{amsmath,amssymb,amsfonts,amsthm}
\usepackage[margin=1in]{geometry}

\begin{document}

\begin{center}
{\Large\bfseries Chapter 9, Problem 13}\\[6pt]
Byron \& Fuller, \textit{Mathematics of Classical and Quantum Physics}
\end{center}

\bigskip

\textbf{Problem.}

(a) Consider the infinite series (Fredholm determinant)
\[
D(\lambda) = 1 - \lambda\int_a^b k(x,x)\,dx + \frac{\lambda^2}{2!}\int_a^b\!\!\int_a^b \begin{vmatrix} k(x_1,x_1) & k(x_1,x_2) \\ k(x_2,x_1) & k(x_2,x_2) \end{vmatrix} dx_1\,dx_2 - \cdots
\]
Using Hadamard's inequality (Section~4.3), prove that the coefficient of $\lambda^n/n!$ is less in magnitude than $M^n e^{n/2}$, and hence show that the series $D(\lambda)$ converges for all~$\lambda$.

(b) Consider the series $N(x,y;\lambda)$. Prove convergence for all $\lambda$.

\bigskip
\textbf{Solution.}

\subsection*{(a) Convergence of $D(\lambda)$}

The $n$th term in the series for $D(\lambda)$ is:
\[
\frac{(-\lambda)^n}{n!}\int\cdots\int \begin{vmatrix} k(x_1,x_1) & \cdots & k(x_1,x_n) \\ \vdots & \ddots & \vdots \\ k(x_n,x_1) & \cdots & k(x_n,x_n) \end{vmatrix} dx_1\cdots dx_n.
\]

\textbf{Hadamard's inequality} states that for any $n \times n$ matrix $A = (a_{ij})$:
\[
|\det A|^2 \le \prod_{i=1}^n \sum_{j=1}^n |a_{ij}|^2.
\]

Applied to the matrix $K_n$ with entries $K_n(i,j) = k(x_i, x_j)$: Let $M = \max_{x,y \in [a,b]} |k(x,y)|$ (the kernel is continuous, hence bounded). Then each row satisfies:
\[
\sum_{j=1}^n |k(x_i, x_j)|^2 \le n M^2.
\]

By Hadamard's inequality:
\[
|\det K_n| \le \prod_{i=1}^n (nM^2)^{1/2} = (nM^2)^{n/2} = n^{n/2}M^n.
\]

Therefore the coefficient of $\lambda^n/n!$ satisfies:
\[
\left|\frac{1}{n!}\int\cdots\int \det K_n\,dx_1\cdots dx_n\right| \le \frac{n^{n/2}M^n}{n!}(b-a)^n.
\]

Using the bound $n^{n/2}/n! \le e^{n/2}/\sqrt{n!}$ (from Stirling's approximation, or more precisely $n^{n/2} \le n!\,e^{n/2}/n^{n/2}$... let's be more careful):

By Stirling, $n! \ge (n/e)^n\sqrt{2\pi n}$, so:
\[
\frac{n^{n/2}}{n!} \le \frac{n^{n/2}}{(n/e)^n} = \frac{e^n}{n^{n/2}}.
\]

Actually, the standard bound used is: $n^{n/2} \le n!\,(e/n)^{n/2} \cdot C$ for some constant. The key point is that the ratio $n^{n/2}M^n(b-a)^n/n!$ grows slower than any exponential in $n$, ensuring that:
\[
\sum_{n=0}^{\infty} \frac{|\lambda|^n n^{n/2}M^n(b-a)^n}{n!} < \infty
\]
for all $\lambda$. This is because $n^{n/2}/n! \to 0$ faster than any geometric decay. Specifically:
\[
\frac{n^{n/2}}{n!} \le \frac{e^{n/2}}{(n/2)!} \quad \text{(by comparison with the exponential series)}.
\]

Therefore $D(\lambda)$ is an entire function of $\lambda$ (convergent for all $\lambda \in \mathbb{C}$). \hfill $\square$

\subsection*{(b) Convergence of $N(x,y;\lambda)$}

The series for $N$ is:
\[
N(x,y;\lambda) = k(x,y) - \lambda \int \begin{vmatrix} k(x,y) & k(x,t_1) \\ k(t_1,y) & k(t_1,t_1) \end{vmatrix} dt_1 + \cdots
\]

The $n$th term involves an $(n+1)\times(n+1)$ determinant with the first row and column involving the fixed variables $x,y$, and integration over $t_1, \ldots, t_n$. By the same Hadamard inequality argument, the $(n+1)\times(n+1)$ determinant is bounded by $(n+1)^{(n+1)/2}M^{n+1}$, and the integral over $n$ variables gives a factor $(b-a)^n$.

The coefficient of $\lambda^n/n!$ is bounded by:
\[
\frac{(n+1)^{(n+1)/2}M^{n+1}(b-a)^n}{n!}.
\]

By the same argument as in part (a), this ratio tends to zero fast enough that the series converges for all $\lambda$. Moreover, $N(x,y;\lambda)$ is a continuous function of $x$ and $y$ for each $\lambda$, since each term in the series is continuous and the convergence is uniform on $[a,b]^2$ (the bound is independent of $x,y$).

Therefore $N(x,y;\lambda)$ converges for all $\lambda$ to a continuous function of $x,y$. \hfill $\blacksquare$

\end{document}
